\section{Cyclotomic interpolation at analytic rank zero}
\label{sec:cy:interpolation}

The cyclotomic argument transfers a nonzero central value across a binary
family. Its input is a uniform upper bound for the normalized valuations
at every nonzero vertex. Its output is nonvanishing at a corank-zero
base. The class used here is Kato's zeta class, and the additional variable
is the real cyclotomic $\Ztwo$-extension of $\Q$.

Retain $T=T_2(E)$ and put $V_2=T\otimes_{\Ztwo}\Qtwo$.
The local weights, finite Sha lengths and period normalization are those
of Section~\ref{sec:preliminaries}, and $D(h)$ is defined in
Equation~\eqref{eq:transfer-normalized-value}. We prove
Theorems~\ref{thm:cyclotomic-bound} and~\ref{thm:cyclotomic-interpolation}
with their stated distinction between the uniform specialization error
and the dimension threshold for a fixed base.

For a two-dimensional representation over $\F_2$, reducibility is
equivalent to being an extension of trivial modules. These are therefore
the two exhaustive residual cases. The proof has three stages: construct
a class in a fixed integral lattice, identify its central determinant,
and apply the binary congruences. Residual reducibility makes the
cyclotomic complex free after localization at $(2)$; the irreducible
case instead uses the bounded clearing factor of
Lemma~\ref{lem:h:pivot-clearing}.

\subsection{The positive cyclotomic complex}

At the coefficient prime $2$, the real place must be included in the
cochain construction. The positive complex defined below removes the
real cochains and retains their connecting invariant line. We will
combine this fixed line with the zeta class to form the determinant
whose central value detects nonvanishing at a corank-zero base.

Put
\[
 \Gamma=\Gal(\Q_\infty/\Q),\qquad
 \Lambda=\Ztwo[[\Gamma]]=\Ztwo[[t]],\qquad
 \mathcal O=\widehat{\Lambda_{(2)}},\qquad
 \Omega=\mathcal O/2\mathcal O=\F_2((t)),
\]
where a generator of $\Gamma$ corresponds to $1+t$.
The completion is $2$-adic. An element of $\mathcal O$ has a finite
negative Laurent part modulo every $2^M$, but need not have a value at
$t=0$. Every central evaluation below is made on a rational Iwasawa
class or on an element of $\Lambda[1/2]$ before extending coefficients
to $\mathcal O$.

For a family in Equation~\eqref{eq:cy:cube}, let
$G=(\Z/2\Z)^b$ have characters indexed by $x\in\F_2^b$, and let
$\psi:G_{\mathbb Q}\to G$ be its universal
quadratic character: evaluation at $x$ is $\chi_{h_x/h_0}$.  Put
\[
 R=\Lambda[G],\qquad
 \mathbb T=T\otimes_{\mathbb Z _2}R,
\]
with the fixed twist $\chi_{h_0}$, the universal twist $\psi$, and the
inverse cyclotomic action.  One may replace $G$ by the actual image of
$\psi$ and subsequently extend scalars.  In a family in
Theorem~\ref{thm:cyclotomic-interpolation}, a nonzero character trivial on that image
already gives the conclusion directly.

Let $S$ be the set of primes dividing $2N h_0\prod_{q\in\mathcal Q}q$.
For $q\in\mathcal Q$, write $g_q\in G$ for its inertia element, so
$x(g_q)=(-1)^{\lambda_q(x)}$. Use continuous cochains with their completed
coefficient topology, and set
\begin{equation}\label{eq:cy:positive-complex}
 C_S^+(\mathbb T)=
 \operatorname{Cone}\bigl(R\Gamma(G_{\mathbb Q,S},\mathbb T)
              \longrightarrow R\Gamma(\mathbb R,\mathbb T)\bigr)[-1].
\end{equation}
The real complex here uses ordinary cochains in nonnegative degrees.
There are no imposed local conditions at the finite places. We write
$C_S(\mathbb T)=R\Gamma(G_{\Q,S},\mathbb T)$ for the unmodified complex.
Throughout this section, \emph{ordinary global cohomology} means this
unmodified cohomology; it imposes no ordinary-reduction hypothesis at~$2$.

\begin{lemma}\label{lem:cy:positive-model}
The complex in Equation~\eqref{eq:cy:positive-complex} is perfect over $R$ and computes derived
coefficient specializations.  There is a class $a\in H^1(C_S^+(\mathbb T))$
whose rational span is the real-place summand.  Its integral index in
that summand depends only on $E$ and the fixed real sign.
\end{lemma}
\begin{proof}
Here are the finiteness and coefficient issues in the use of global
cochains.  For a finite coefficient quotient of $R$, continuous cochains
are filtered unions of finite products of free coefficient modules.
Refinement of finite partitions makes these modules flat; over an Artin
local coefficient ring they are free.  For finite residual coefficients,
global cohomology is finite in each degree, and restriction to the real
places is an isomorphism in degrees greater than two, by the real-place
clause of global duality \cite[\S5.7.1.7]{Nekovar2006}.
Thus the residual fiber has bounded
amplitude and finite-dimensional cohomology.  The same amplitude bound
holds for every finite coefficient module, by a composition series and
the cohomology long exact sequence.

This gives a bounded finite free model as follows.  Over the residue
field split off the acyclic disks, leaving its finitely many cohomology
spaces.  Lift each invertible disk differential over an Artin coefficient
quotient and cancel it.  Invertibility of a lifted inverse follows from
nilpotence of the maximal ideal.  Equivalently one can perform this on a
free resolution; a minimal free model cannot have an additional term
whose reduction has zero cohomology, because its residual differential
is zero.  Hence the remaining ranks are the finite residual ranks and
are bounded independently of the coefficient quotient.  Minimal models
identified by a homotopy equivalence over an Artin local ring are
isomorphic.  Lifting their changes of basis therefore makes these models
compatible along $R/\mathfrak m^n$.  Taking the inverse limit gives a
bounded finite free model over~$R$.  The cochain construction with finite
free coefficients and the same cancellation also proves derived base
change.  This is a fixed-$S$ argument; it does not assert uniform bounds
on the number of generators as $S$ varies.

All the characters of the family have the same real sign, so the action
at the real place is the fixed action on $T\chi_{h_0}$ extended to $R$.
Choose a nonzero vector in its integral invariant line.  Its connecting
image in the fiber is $a$.  There are no global invariant vectors in the
rational elliptic representation, and consequently this connecting map
is injective rationally.  The choice and its index are fixed.  Notice also
that $H^1(\mathbb R,\mathbb T)$ is killed by~$2$; thus twice any ordinary
cohomology class lifts to the positive complex.  The real-place global duality input is its corrected form at~$2$
\cite[\S\S5.7.1--5.7.5]{Nekovar2006}; the invariant line and the lifting
assertion follow from the long exact sequence of this ordinary-real cone.
Complete real Tate cochains, used in the duality formulation, become
acyclic after inverting~$2$ and would not retain this invariant line.
\end{proof}

\subsection{Fixed-form zeta classes and the integral coefficient map}

To control denominators as the tame support grows, we construct every
class in the cohomological lattice of one fixed modular form.  Let $f$ be the fixed weight-two form of $E$, and choose
its fixed integral cohomological lattice, with the Tate twist making its
rational representation~$V_2$.  Kato's classes (8.1.2)--(8.1.3), for a fixed
symbol $\xi\in\mathrm{SL}_2(\mathbb Z)$, exist for arbitrary $m$ and any
$S$ containing the prime factors of $mN$.  The trace in Section~8.9 is
a homomorphism into the lattice
\[
 \mathbb Z _2[\operatorname{Gal}(\mathbb Q(\mu_m)/\mathbb Q)]
       \otimes V_{2,\mathbb Z _2}(Y_1(N))(1).
\]
Section~8.11 then projects to the lattice of the fixed form.
Proposition~8.12 supplies its norm relations.  These assertions hold at
the coefficient prime~$2$ \cite[Sections~5.1, 8.1, 8.9--8.12]{Kato2004}.
In particular, the auxiliary modular-curve level used to define a trace
does not become the level of the projected newform.

Choose an isomorphism of the resulting fixed rational representation with
$V_2$.  A single power $2^{c_E}$ sends its integral lattice into~$T$.
Choose fixed symbols with nonzero projections on the required period
lines; their existence and the construction with both signs are made
explicit in \cite[Sections~13.6 and 13.9]{Kato2004}.  Taking a fixed linear
combination, if needed, costs another fixed denominator.

The following integral coefficient map encodes the binary family. Choose
$m=|h_0|\prod_{q\in\mathcal Q}q$, which is odd and divisible by the
conductors of all the tame characters in the family, including
$\chi_{h_0}$.  At a sufficiently large
cyclotomic level write $A_{m,n}=\Gal(\Q(\mu_{m2^n})/\Q)$
and let $\Gamma_n$ be its quotient in the real cyclotomic tower.  Shapiro identifies the cohomology of the cyclotomic field
with global cohomology with induced group-ring coefficients.  Push the
class along the ring homomorphism
\begin{equation}\label{eq:cy:coefficient-map}
 \mathbb Z _2[A_{m,n}]\longrightarrow
       \mathbb Z _2[\Gamma_n\times G],\qquad
 [\sigma]\longmapsto
       \chi_{h_0}(\sigma)[\sigma|_{\Gamma_n}][\psi(\sigma)].
\end{equation}
Use inverse actions on both sides if this is required by the chosen
Shapiro convention.  This is a homomorphism of integral coefficient
modules and acts on actual cohomology classes.  It requires neither
division by $|G|$ nor surjectivity of a cohomological specialization map.
The prime $2$ belongs to the fixed omitted support, so Proposition~8.12
gives exact norm compatibility as $n$ increases.  The pushed classes
therefore give an integral smoothed Iwasawa class with coefficients in
$\mathbb T$, up to the single fixed power~$2^{c_E}$.

We have placed all character constituents in the lattice of the same
modular form. The next normalization issue is smoothing: its
removal must be harmless both after localization at $(2)$ and at the
central value $t=0$. These are distinct requirements, checked next.

We next remove smoothing.  In weight two and with $r=r'=1$, its factors
are
\begin{equation}\label{eq:cy:smoothing}
 (c^2-c\sigma_c)(d^2-d\sigma_d).
\end{equation}
This is the specialization of
\cite[Theorem~6.6 and (4.2.4)]{Kato2004}.  Choose a fixed
$e\geq\max(2,v_2(N))$, increasing it to accommodate the finitely many
fixed symbol levels.  The Chinese remainder theorem permits $c,d>1$,
prime to $6Nm$, such that
\[
 c,d\equiv1\pmod{3N_{\mathrm{odd}}m},\qquad
 c,d\equiv1+2^e\pmod{2^{e+1}}.
\]
Their tame actions in \eqref{eq:cy:coefficient-map} are trivial.  If
$\sigma_c$ acts in $\Gamma$ as $(1+t)^{\kappa(c)}$, then
$v_2(\kappa(c))=e-2$.  Modulo the maximal ideal of $\mathcal O[G]$ its
smoothing factor is
\[
 1-(1+t)^{\kappa(c)}\ne0\quad\hbox{in }\mathbb F _2((t)).
\]
It is therefore a unit in $\mathcal O[G]$.  Its central value is
$c(c-1)$, of valuation exactly~$e$.  The same statements hold for $d$
and for inverse-translate conventions.  Thus unsmoothing introduces
neither a $2$-adic denominator at $(2)$ nor a growing central denominator.
The integers $c,d$ vary with $m$; the congruences fix the needed valuations,
not their exact cyclotomic group elements.

Denote the resulting rational ordinary class by $z_S$.  \citet[Theorem~9.7]{Kato2004} applies to $k=2,r=r'=1$.  In the character constituent
$x$, its central reciprocity formula reads
\begin{equation}\label{eq:cy:reciprocity-center}
 v_2\bigl(\exp_2^*z_{S,x}(0)\bigr)
 =v_2(\mathcal L(h_x))+
  \sum_{\substack{q\in\mathcal Q\,,\ q\nmid h_x}}
       v_2\bigl(\#E^{(h_x)}(\mathbb F _q)\bigr)+O_E(1)
\end{equation}
when the value is nonzero, and a zero central value gives a zero dual
exponential.  The omitted primes dividing $2N$ are included in the error.
To check the normalization, the character sum in
\eqref{eq:cy:coefficient-map} is the unnormalized sum, as in Kato's trace.
Its tame quadratic Gauss sum supplies $\sqrt{|h_x|}$ up to a root of
unity.  Tame conductors and their Gauss sums are $2$-adic units.  The
period vector is the fixed symbol vector of $f$.  More explicitly,
in the required sign line write the period represented by
$\delta(f,1,\xi)^{\pm}$ as $c_\xi^{\pm}\Omega_E^{\pm}$.
The nonzero scalar $c_\xi^{\pm}$ is fixed: the character-sum formula
of \cite[Theorem~6.6]{Kato2004} uses this same $\delta(f,1,\xi)^{\pm}$
for every tame character.  The tame character enters through its
$L$-value and Gauss sum, not through a newly chosen period vector.
The comparison with the minimal differential of $E^{(h_x)}$ at $2$ has bounded valuation,
because there are only finitely many local twisting classes there.
At a good unused prime the omitted factor is
\[
 1-\frac{a_q(E)\chi_{h_x}(\operatorname{Fr}_q)}q+\frac1q
       =\frac{\#E^{(h_x)}(\mathbb F _q)}q;
\]
at a good ramified prime the local $L$-factor is~$1$.  These account for
every varying factor in \eqref{eq:cy:reciprocity-center}.

\subsection{Rational divisibility and the local Euler corrections}

We now exclude poles of the zeta determinant away from $(2)$ and
account for every omitted Euler factor.  Write $A=\Lambda[1/2]$.  Theorem~12.5(3) of
\cite{Kato2004} is a height-one divisibility over~$A$; its possible local
$H^2$ correction comes from the arithmetic place~$2$. It need not be
supported at the Iwasawa ideal~$(2)$. The global groups
there use $j_*$ outside that prime.  The stronger integral assertion of
Theorem~12.5(4) is not used.  For a non-CM elliptic curve the characteristic-zero
image condition is supplied by the open image theorem, independently of
the residual image.  For CM the required rational statement at every
prime, including~$2$, is supplied by
\cite[Theorem~2.6 and Remark~2.7]{BurungaleTian2026}.  Neither statement
asserts uniform integrality in a quadratic group ring. The rational
cohomology statements, including \citet[Theorem~12.4]{Kato2004}, give rank one in degree one and torsion in degree
two, with no other positive cohomological rank. The real-place
modification adds its invariant line. Hence every generic character
constituent of the positive complex has rank two in degree one and no
other positive rank.

Fix a character constituent $h$ and write $V_h=V_2\chi_h$.  Let $C_{j_*}$
be the rational Iwasawa complex in Kato's formulation and $C_S$ the
rational full global complex.  At an odd prime put
\[
 W_q=H^1(I_q,V_h).
\]
The localization triangle has local cone
\begin{equation}\label{eq:cy:local-block}
 Q_q=R\Gamma\bigl(G_{\mathbb F _q},
                   W_q\otimes A(\Gamma^\iota)\bigr)[-1]
   =\left[W_q\otimes A
       \xrightarrow{\ 1-\gamma_q\operatorname{Fr}_{W_q}\ }
             W_q\otimes A\right]
\end{equation}
in degrees one and two.  Here $\gamma_q$ denotes the scalar cyclotomic
Frobenius in the chosen convention.  Its exponent is nonzero.  Consequently
$D_q=\det(1-\gamma_q\operatorname{Fr}_{W_q})$ is a nonzero element of~$A$,
and the map in \eqref{eq:cy:local-block} is injective.  This remains true
after localization at a height-one prime dividing $D_q$: localization
does not mean reduction modulo that prime.  Thus
\[
 H^1(Q_q)=0,\qquad
 \operatorname{length}_{A_{\mathfrak p}}H^2(Q_q)_{\mathfrak p}
       =v_{\mathfrak p}(D_q).
\]
The long exact sequence gives
\begin{equation}\label{eq:cy:local-sequence}
 \begin{gathered}
 H^1(C_{j_*})\simeq H^1(C_S),\\
 0\longrightarrow H^2(C_{j_*})\longrightarrow H^2(C_S)
       \longrightarrow\bigoplus_{q\in S\setminus\{2\}}H^2(Q_q)
       \longrightarrow H^3(C_{j_*}).
 \end{gathered}
\end{equation}
In particular, the length of $H^2(C_S)$ is at most the length of
$H^2(C_{j_*})$ plus the sum of the local lengths.  No assertion that
$H^3(C_{j_*})$ vanishes is required here.

At a good unused prime, $W_q=V_h(-1)$, so, on including the quadratic
Frobenius sign in $Z$,
\begin{equation}\label{eq:cy:local-euler}
 D_q=P_q(Z),\qquad P_q(Z)=1-\frac{a_q(E)}q Z+\frac1q Z^2.
\end{equation}
At a good active prime, inertia acts by $-1$ on $V_h$.  Its rational
inertia cohomology is zero, so $W_q=0$ and $D_q=1$.

At fixed bad primes the general expression is $H^1(I_q,V_h)$ in
\eqref{eq:cy:local-block}.  Replacing it indiscriminately by $V_h^{I_q}(-1)$
would be incorrect.  For example, at split multiplicative reduction,
\[
 V_h^{I_q}=\mathbb Q _2(1),\qquad
 H^1(I_q,V_h)=\mathbb Q _2(-1),
\]
and the correct determinant is $1-\gamma_q/q$, with nonzero central
value.  The incorrect replacement would give $1-\gamma_q$.
In general $D_q(0)\ne0$ follows from local duality and the Weil pairing:
$H^2(\mathbb Q _q,V_h)=H^0(\mathbb Q _q,V_h)^*=0$, since elliptic torsion
over a local field is finite.  The fixed bad-prime corrections therefore
have nonzero central values for every reduction type.

Let $z_{\mathrm{BK},h}$ denote a generator of the rational Kato zeta
line in its real cyclotomic component.  The rational comparison with
the fixed-form class constructed above is made only after taking the
character constituent.  Both global Iwasawa $H^1$ modules have rank one.
Their reciprocity laws at cyclotomic finite-order characters agree up to
the omitted Euler factors, a nonzero constant, and a group-like unit.
This follows by coprime Gauss-sum factorization and the fixed period
comparison; restricting conductor exponents to either fixed parity
removes any factor $\chi_h(2)^n$.  Infinitely many of these tests have
nonzero central value by cyclotomic nonvanishing, as used in
\cite[Theorem~13.5(2)]{Kato2004}.  After clearing the finitely many
denominators of the scalar ratio, the difference is a power series
vanishing at infinitely many finite-order characters.  Weierstrass
preparation forces that difference to be zero.  This proves the generic
comparison without identifying the integral lattices of varying
twisted newforms.  Its nonzero scalar in $\mathbb Q _2$ may depend on
$h$ and need not have a uniformly bounded $2$-valuation: every such
scalar is a unit of $A$ and does not affect height-one divisibility
over~$A$.  The uniform integral and central bounds instead concern
the actual fixed-form class and the fixed period vector in
\eqref{eq:cy:reciprocity-center}.

Equations~\eqref{eq:cy:local-sequence}--\eqref{eq:cy:local-euler} now show exactly
how the imprimitive factors enter divisibility.  At any height-one
$\mathfrak p$ of $A$, enlarging the global complex increases its
degree-two length by at most $\sum_qv_{\mathfrak p}(D_q)$, whereas
multiplying the rank-one zeta generator by $\prod_qD_q$ adds exactly
that quantity to its index.  To cover Kato's coefficient-prime correction, multiply the
class also by a fixed element $F_2\in\Lambda$ divisible over $A$ by the
characteristic ideal of the relevant local $H^2_{\mathrm{Iw}}$.
It may be chosen with $F_2(0)\ne0$.  Indeed derived local control gives
\[
 H^2_{\mathrm{Iw}}(\mathbb Q _2,T\chi_h)_\Gamma\otimes\mathbb Q _2
       =H^2(\mathbb Q _2,V_h)=0,
\]
so the characteristic ideal has no factor $t$.  Local twisting classes
at $2$ are fixed, and one element covers their finite set.

Likewise any additional fixed Euler factors in the rational comparison
can be cleared by a fixed product with nonzero central value, using the
actual blocks \eqref{eq:cy:local-block}.  Denote the total fixed multiplier
by~$F_E$.  Multiplying by it only changes the fixed constants in
\eqref{eq:cy:reciprocity-center} and the integrality bound.  For the corrected
class, still written $z_S$, we have at every height-one $\mathfrak p$ of~$A$
\begin{equation}\label{eq:cy:divisibility}
 \operatorname{length}_{A_{\mathfrak p}}H^2(C_S)_{\mathfrak p}
 \leq
 \operatorname{length}_{A_{\mathfrak p}}
       \bigl(H^1(C_S)/Az_{S,h}\bigr)_{\mathfrak p}.
\end{equation}

There is an integral convention check for the varying unused primes.
In $\mathcal O[G]$, both $P_q(Z)$ and $Z^2P_q(Z^{-1})$ are units: their
reductions are nonzero polynomials evaluated at the nonconstant
cyclotomic Frobenius.  The identity
\[
 P_q(Z)-Z^2P_q(Z^{-1})=\frac{q-1}{q}(1-Z^2)
\]
shows that their ratio $r_q$ is $1$ modulo~$2$.  If $g_q$ is the universal
inertia element, then
\begin{equation}\label{eq:cy:reciprocal-correction}
 U_q=1+\frac{r_q-1}{2}(1+g_q)\in\mathcal O[G]^\times
\end{equation}
is integral.  At an active character it is $1$; at an unused character
it is $r_q$.  Its unused central value is $1$, since $Z=\pm1$ there.
Thus inverse Frobenius conventions can be reconciled before character
evaluation, without a denominator for each prime.  Products of these
units have the same property.  The resulting rational character
coordinates have the Euler convention in \eqref{eq:cy:divisibility}.

\begin{lemma}[Fixed-lattice zeta class]\label{lem:cy:kato-class}
For every binary family and its ramification set $S$, the corrected
class $z_S$ satisfies the central reciprocity formula \eqref{eq:cy:reciprocity-center} and the rational divisibility
\eqref{eq:cy:divisibility}. For some $c_E\ge0$ depending only on $E$,
\[
 2^{c_E}z_S\in
 H^1\!\left(C_S(\mathbb T)\otimes_R^{\mathbf L}\mathcal O[G]\right).
\]
The same exponent works for all $S$, binary families, and fixed factors
$h_0$ with the specified local classes. Here the displayed cohomology
is derived scalar extension of the Iwasawa complex.
\end{lemma}
\begin{proof}
The integral coefficient map uses the fixed-form lattice and no
character averages. Unsmoothing and every varying-prime convention
correction are units over $\mathcal O[G]$. The multiplier $F_E$ is fixed.
The preceding construction therefore gives the displayed assertion
with a fixed exponent, as well as the two stated comparison formulas.
\end{proof}

\subsection{The determinant and its central specialization}

The zeta class now has two controls: rational divisibility will exclude
poles of its determinant coordinate away from $(2)$, and the fixed
lattice bounds the class's integral denominator at $(2)$. We next
identify its determinant coordinate at
$t=0$. This is the step that converts those cohomological controls
into a statement about $D(h)$.

Choose a basis of the invertible $R$-module
$\det_R^{-1}C_S^+(\mathbb T)$; $R$ is local, so such a basis exists.
Lift $2z_S$ to rational positive cohomology and take its wedge with $a$.
The factor $2$ is fixed, and two lifts differ by a real-place class,
which does not change this wedge.  Let $u$ be its generic determinant
coordinate, including the fixed multiplier $F_E$ above.  The real-place
summand introduces no height-one torsion over~$A$.  Thus
\eqref{eq:cy:divisibility} says that $u_x$ has no pole at any height-one
prime of $\Lambda$ except possibly $(2)$.  Since $\Lambda$ is normal,
\begin{equation}\label{eq:cy:rational-integrality}
 u_x\in\Lambda[1/2]\qquad\hbox{for every }x.
\end{equation}
The two cycles have a common denominator bounded by a fixed power
$2^{C_{\mathrm{int}}}$ after extension to $\mathcal O[G]$. This follows
from Lemma~\ref{lem:cy:kato-class} and the factor $2$ needed for the lift.

\begin{proposition}[Central determinant formula]\label{prop:cy:specialization}
There is a constant $C_{\mathrm{loc}}$, depending only on the fixed curve
and local classes, with the following property. Suppose
$c_2(E^{(h_x)})=0$. If $L(E^{(h_x)},1)=0$, then $u_x(0)=0$; if
$L(E^{(h_x)},1)\ne0$, then
\begin{equation}\label{eq:cy:determinant-center}
 \bigl|v_2(u_x(0))-D(h_x)\bigr|\leq C_{\mathrm{loc}}.
\end{equation}
The rational positive complex at $t=0$ has cohomology only in degree one,
of dimension two, so these are specializations of the same determinant.
\end{proposition}
\begin{proof}
Abbreviate $T_h=T_2E^{(h_x)}$ and $A_h=E^{(h_x)}[2^\infty]$.
Use the integral Kummer subgroup
\[
 K_v=\operatorname{im}\bigl(E^{(h_x)}(\mathbb Q _v)^{\wedge}_2
                              \longrightarrow H^1(\mathbb Q _v,T_h)\bigr)
\]
at every finite place.  For $v\ne2$ it is all of $H^1(\mathbb Q _v,T_h)$.
One way to see this integrally is the inverse-limit Kummer exact sequence:
its quotient is $T_2H^1(\mathbb Q _v,E^{(h_x)})$, which is zero because
local Tate duality identifies the $2$-primary local Weil--Ch\^atelet group
with the dual of the finite $2$-completion of the local point group.
At $2$ the quotient
\[
 L_2=H^1(\mathbb Q _2,T_h)/K_2
\]
is a free rank-one $\mathbb Z _2$-module.  The same Kummer sequence and
local Tate duality identify it with the $\mathbb Z _2$-dual of the free
part of $E^{(h_x)}(\mathbb Q _2)^{\wedge}_2$.

The dual exponential identifies $L_2\otimes\mathbb Q _2$ with the
one-dimensional de Rham target. To specify its integral comparison, put
\[
 B_2=E^{(h_x)}(\mathbb Q _2)^{\wedge}_2/\mathrm{tors}
\]
and choose a minimal differential $\omega_h$.
The Bloch--Kato exponential agrees with the Kummer image of the classical
elliptic exponential, which is the inverse of the formal logarithm on a
sufficiently small open subgroup. This comparison holds for every
reduction type \cite[Definition~3.10 and Example~3.11]{BlochKato1990}.
The dual exponential is the adjoint of that exponential under local
Tate duality \cite[\S II and Proposition~II.5]{Berger2003}.
Extend the formal logarithm to the rationalized local point group and write
$\log_{\omega_h}(B_2)=\lambda_h\mathbb Z _2$.
The integral identification
$L_2=\operatorname{Hom}_{\mathbb Z _2}(B_2,\mathbb Z _2)$ then gives
\[
 \exp_2^*(L_2)=\lambda_h^{-1}\mathbb Z _2\omega_h,
\]
using the Weil-pairing identification of the de Rham target with
$\mathbb Q _2\omega_h$. Thus the relative lattice valuation is fixed
by the local curve and differential. Since only finitely many local
twist classes occur, these valuations are uniformly bounded.

The assumption on the Selmer corank implies that the compact finite
Selmer group is torsion and that the discrete finite Selmer group has
length $s(h_x)$.  Rational Poitou--Tate consequently identifies ordinary
global $H^1$ with the singular line at $2$, and gives ordinary global
$H^2=0$ rationally.  Adding the real invariant line gives the claimed
dimension two in degree one.  This equals the generic rank, so derived
base change of the perfect complex specializes its determinant without
a change of rational cohomological rank.

Here is the integral length calculation.  Let $J$ be the image of the
free part of global $H^1(G_{\mathbb Q,S},T_h)$ in $L_2$, and put
$i=\operatorname{length}(L_2/J)$.  The compact/discrete Poitou--Tate
sequence for the mutually orthogonal Kummer conditions contains
\[
\begin{aligned}
 0\longrightarrow L_2/J
 &\longrightarrow \operatorname{Sel}_{2^\infty}(E^{(h_x)}/\mathbb Q)^\vee
 \longrightarrow H^2(G_{\mathbb Q,S},T_h)\\
 &\longrightarrow \bigoplus_{q\in S}H^2(\mathbb Q _q,T_h)
 \longrightarrow H^0(\mathbb Q,A_h)^\vee\longrightarrow0,
\end{aligned}
\]
with the usual real Tate-cohomology terms inserted.  Those terms are
killed by $2$ and have dimensions bounded by the fixed dimension of~$T$.
Taking lengths and using local duality gives
\begin{equation}\label{eq:cy:pt-length}
 \operatorname{length}H^2(G_{\mathbb Q,S},T_h)
   =s(h_x)-i+\sum_{q\in S}\operatorname{length}H^0(\mathbb Q _q,A_h)
      +O_E(1).
\end{equation}
This sequence follows from integral Poitou--Tate duality with the finite
Kummer conditions and the real Tate-cohomology correction
\cite[\S\S5.7.1.6, 5.7.4--5.7.5 and 6.9.1]{Nekovar2006}.
The free real invariant line is already accounted for by $a$; only the
bounded real torsion terms enter the error above.
No unramified condition replaces
the Kummer condition at a bad prime.  Global torsion is bounded by its
injection into $E^{(h_x)}(\mathbb Q _2)[2^\infty]$, whose order is bounded
by the fixed local twisting classes.  The finite terms at $2$ and at
fixed bad primes are bounded for the same reason.  Passing between
ordinary and positive cohomology changes only these bounded real and
global torsion terms.

Suppose first that the central value is nonzero.  The coordinate of
$z_{S,x}(0)$ in the free global line has valuation
\[
 v_2(\exp_2^*z_{S,x}(0))-i+O_E(1).
\]
For the inverse determinant, torsion in degree two subtracts its length
and torsion in degree one adds its length.  This sign can be checked on
the complex $[\mathbb Z _2\xrightarrow{2^a}\mathbb Z _2]$ in degrees one
and two: its inverse determinant lattice maps to $2^a\mathbb Z _2$ in
the rational determinant.  A vector's coordinate thus loses~$a$.
Using \eqref{eq:cy:pt-length}, the two appearances of $i$ cancel, leaving
\[
 v_2(u_x(0))=
 v_2(\exp_2^*z_{S,x}(0))-s(h_x)
       -\sum_{q\in S}\operatorname{length}H^0(\mathbb Q _q,A_h)+O_E(1).
\]

At an active good prime, quadratic inertia acts by $-1$ on the Tate
module.  Its invariants on $A_h$ are therefore exactly $E[2]$, and taking
Frobenius invariants gives length~$t_q$.  At an unused good prime, reduction
identifies the local $2$-primary torsion with
$E^{(h_x)}(\mathbb F _q)[2^\infty]$, of length
$v_2(\#E^{(h_x)}(\mathbb F _q))$.  These unused terms cancel exactly
those in \eqref{eq:cy:reciprocity-center}.  The active terms, including the primes
of $h_0$, sum to $2w(h_x)$.  All other terms have already been bounded
independently of the growing support.  This proves
\eqref{eq:cy:determinant-center}.

If the central value is zero, reciprocity makes the dual exponential
zero.  The rational ordinary specialization then satisfies the finite
condition at $2$, and it satisfies the finite conditions elsewhere
automatically.  It lies in the rational finite Selmer group, which is
zero.  Its wedge with the real-place class is therefore zero.  The
constant-rank determinant specialization proves $u_x(0)=0$.
\end{proof}

\subsection{Residual freeness and the lower bound}

The central formula expresses $D(h)$ through the determinant valuation.
When the residual representation is an extension of trivial modules,
we can bound that valuation directly: the positive complex becomes a
free rank-two module in degree one after extension to $\mathcal O[G]$.
The fixed-lattice class then supplies the lower bound without any loss
from a growing ramification support.

\begin{lemma}[Residual concentration]\label{lem:cy:residual-free}
If $E[2]$ is an extension of two trivial $\F_2$-modules, then there is an
isomorphism in the derived category
\[
 C_S^+(\mathbb T)\otimes_R\mathcal O[G]
       \simeq \mathcal O[G]^2[-1].
\]
\end{lemma}
\begin{proof}
The ring $\mathcal O[G]$ is local, with residue field $\Omega$ obtained
by setting $2=0$ and every $g=1$.  Thus the quadratic characters disappear.
Since $E[2]$ has a filtration with trivial factors, it is enough first to
compute the positive residual complex for the trivial representation.

Let $K_n$ be the real cyclotomic layer of degree $d_n=2^n$.  The number
of finite primes of $K_n$ above a fixed $S$ is bounded with $n$: at $2$
there is total ramification, and the Frobenius of each fixed odd prime
has nonzero image in $\Gamma$.  The $2$-ranks of the class groups of
$K_n$ are also bounded. Let $X$ be the inverse limit, under norm maps, of
their $2$-primary class groups. The unique prime above~$2$ is totally
ramified in this tower, so class-field-theoretic control gives
\[
 X_{\Gamma^{2^n}}\simeq\Cl(K_n)[2^\infty]
\]
\cite[Proposition~3.2.11 and Remark~3.2.12]{SharifiIwasawa}.
Ferrero--Washington gives $\mu(X)=0$ for this abelian tower
\cite{FerreroWashington1979}. The Iwasawa structure theorem then makes
$X$ finitely generated over $\mathbb Z _2$, so all its coinvariants have
uniformly bounded $2$-rank. Passing to the quotient $S$-class groups
preserves this upper bound.

The signature map of units in $K_n$ is surjective.  For $n\geq1$, let
$\zeta$ be a primitive $2^{n+2}$th root and put
$\epsilon=1+\zeta+\zeta^{-1}$.  If $f_1(X)=X^2-2$ and
$f_{n+1}(X)=f_n(X^2-2)$, then $f_n$ is the minimal polynomial of
$\zeta+\zeta^{-1}$ and $f_n(-1)=-1$.  Hence
$N_{K_n/\mathbb Q}(\epsilon)=-1$.  Its signature vector has odd
augmentation.  The embeddings form a regular set for the cyclic
$2$-group $\operatorname{Gal}(K_n/\mathbb Q)$, whose group algebra over
$\mathbb F _2$ is local.  A vector with odd augmentation is a unit in
this group algebra, so its translates span every signature.  For $n=0$
the assertion follows from the unit~$-1$.

The signature calculation supplies the real-place surjectivity needed
for the positive complex. The class-group and support terms are bounded
along the tower. We can therefore identify the one cohomological degree
whose dimension grows with the layer degree, and then recover the rank
after inverting $t$.

Kummer theory now determines the cohomological growth.
Writing $s_n$ for the number of finite support places and $c_n$ for the
$2$-rank of the $S$-class group, one has
\[
 \dim H^1(G_{K_n,S},\mu_2)=d_n+s_n+c_n.
\]
The real degree-one restriction is surjective by the signature result.
In degree two, the Kummer sequence has class-group term of dimension
$c_n$ and Brauer term.  The Brauer invariants at the real and finite
support places have their single sum relation.  Since a finite support
place above $2$ is present, restriction of the Brauer term to the real
places is surjective and has kernel of dimension $s_n-1$.  Thus the
degree-two kernel of restriction to infinity is bounded.  Restriction
is an isomorphism in higher degrees by global duality.  The degree-zero
restriction is injective.  The cohomology of the positive fiber is
therefore bounded in every degree except one; its degree-one dimension
is $d_n+O_S(1)$.

Apply Lemma~\ref{lem:cy:positive-model} modulo $2$.  Derived specialization to the
$n$th layer is specialization modulo $t^{2^n}$.  For a perfect complex
over $\mathbb F _2[[t]]$, the universal-coefficient exact sequence shows
that a rank-$r$ cohomology module contributes $r2^n+O(1)$ to these
dimensions; the torsion terms contribute only $O(1)$.  It follows that
after inverting $t$ the trivial-constituent fiber has dimension one in
degree one and vanishes in every other degree.  Extensions preserve
this concentration, so the fiber for $E[2]$ has dimension two in degree
one.  A minimal finite free complex over the local ring $\mathcal O[G]$
has zero residual differentials.  The concentration just proved forces
all its terms except the rank-two term in degree one to vanish.
No denominator is introduced by this exact freeness statement.
\end{proof}

\begin{proof}[Proof of Theorem~\ref{thm:cyclotomic-bound}]
Apply Lemma~\ref{lem:cy:residual-free} to the single discriminant $h$.
The fixed-lattice construction and the free complex give
$2^{C_{\mathrm{int}}}u_h\in\mathcal O$ with $C_{\mathrm{int}}$ depending
only on $E$. Equation~\eqref{eq:cy:rational-integrality} also gives
$u_h\in\Lambda[1/2]$. Since
$\mathcal O\cap\Lambda[1/2]=\Lambda$, it follows that
$v(u_h(0))\ge-C_{\mathrm{int}}$. Analytic rank zero supplies Selmer
corank zero by the forward theorem of Gross--Zagier--Kolyvagin.
Proposition~\ref{prop:cy:specialization} now gives
$D(h)\ge-C_{\mathrm{int}}-C_{\mathrm{loc}}$.
All constants come from the fixed form, the fixed real period line,
or the finitely many local twisting classes at $2N$.
\end{proof}

\subsection{Transfer to the base}

It remains to prove the missing-vertex theorem. Residual freeness gives
a single integral family in the reducible case. In the irreducible
case we instead need a bounded clearing factor after eliminating the
new local blocks. The bound on $v_2(q^2-1)$ enters precisely here: the
next lemma controls the negative Laurent powers in those inverse
blocks at each fixed precision.

\begin{lemma}\label{lem:cy:cyclotomic-poles}
Suppose that $q$ is an added prime and $v_2(q^2-1)\le L$.  At
augmentation $G\to1$, the inverse singular Frobenius block contributed
by $q$ has, modulo each fixed $2^M$, negative Laurent support bounded
in terms of $L$ and $M$.  This bound is independent of $q$.
\end{lemma}

\begin{proof}
Identify $\Gamma$ with $1+4\Ztwo$ using a fixed generator.  If
$\gamma_q=(1+t)^{\alpha_q}$ is the image of arithmetic Frobenius at
$q$, then
\[
 v_2(\alpha_q)=v_2(q^2-1)-3
\]
when the generator is $5$; changing the generator multiplies
$\alpha_q$ by a unit.  In particular $\alpha_q\ne0$.  The singular
quotient is represented by a block of bounded size whose determinant
is the Euler polynomial in Equation~\eqref{eq:cy:local-euler}.
Modulo $2$, with $z=(1+t)^{\alpha_q}$, it is
\[
 1+\overline{a_q}z+z^2.
\]
If $\overline{a_q}=1$, its constant coefficient is $1$.  If
$\overline{a_q}=0$, it is $(1+z)^2$, whose $t$-order is
$2^{v_2(\alpha_q)+1}\le2^{L-2}$.  Thus the inverse determinant has
uniformly bounded pole order modulo $2$.  Lift its inverse modulo
$2^M$ by a geometric expansion of length at most $M$.  Each term
contains only a bounded number of inverse residual factors, and all
numerators have nonnegative powers of $t$.  The adjugate formula gives
the same bound for every entry of the inverse block.
\end{proof}

\begin{proof}[Proof of Theorem~\ref{thm:cyclotomic-interpolation}]
At every nonzero vertex, analytic rank zero gives
$c_2(E^{(h_x)})=0$ by the forward rank-zero theorem. Thus
Proposition~\ref{prop:cy:specialization} applies at every vertex.

First assume $E[2]$ is an extension of trivial modules.
Lemma~\ref{lem:cy:residual-free} and the fixed-lattice construction give
$2^{C_{\mathrm{int}}}u\in\mathcal O[G]$.  Moreover, the character coordinates of $u$
are in $\Lambda[1/2]$.  Fourier inversion, performed rationally, shows
$u\in\Lambda[1/2][G]$; intersecting coefficientwise with
$\mathcal O[G]$ gives
\begin{equation}\label{eq:cy:universal-integrality}
 2^{C_{\mathrm{int}}}u\in\Lambda[G].
\end{equation}
The coefficients at $t=0$ therefore obey the bounded-degree Boolean
congruences of Lemma~\ref{lem:h:boolean-coefficients}. For any fixed
precision and sufficiently large $b$, there is a nonzero $x$ for which
$u_x(0)$ agrees with $u_0(0)$ to that precision.  The fixed multiplier
in Equation~\eqref{eq:cy:universal-integrality} affects the required
precision but has no dependence on $b$.

Now suppose $E[2]$ is irreducible.  At augmentation, represent the
positive complex as the old complex, with ramification set excluding
$\mathcal Q$, extended by the direct sum of singular blocks at the
new primes.  The differential is upper triangular with the old
complex as subcomplex. The finite-model argument of
Lemma~\ref{lem:h:finite-models} lifts this model over $R$: stabilize the
finite free models and lift their graded changes of basis. These lifts
are invertible because the augmentation ideal lies in the Jacobson
radical. Cancel all additional lifted unit disks before retaining the
old-plus-singular grading; their augmented cross terms are zero, so the
remaining augmentation is exactly the prescribed block model. Thus no
stabilization disk enlarges the old complex.
Lemma~\ref{lem:h:bounded-poles} then applies to this block model.  All the new
blocks are invertible over $\mathcal O[G]$: modulo its maximal ideal
their Frobenius determinants are nonzero in $\Omega$.
Schur elimination leaves a complex of bounded size. Its augmentation
is the fixed old complex, and its entries have bounded negative
Laurent support modulo each fixed $2^M$. Under the determinant-line
identification, write
\[
 u=\eta u_{\mathrm{bd}},\qquad \eta\in\mathcal O[G]^\times,
\]
where $u_{\mathrm{bd}}$ is the coordinate on this bounded complex and
$\eta$ is the determinant factor of the eliminated blocks. No uniform
Laurent bound on $\eta$ is needed.

The pole bound remains uniform as primes are added.  At augmentation the inverse new block is a direct
sum of the inverses in Lemma~\ref{lem:cy:cyclotomic-poles}.  At any
character its perturbation has entries divisible by $2$ and no
negative powers of $t$.  Modulo $2^M$, the inverse uses a geometric
series of length at most $M$, independently of the number of blocks.
Each product therefore uses a bounded number of inverse entries.
Sums of an arbitrary number of such terms cannot decrease the
nonarchimedean valuation or the common lower Laurent bound.  This is
the required bound for Schur elimination.

At $t=0$ the old positive complex has rational cohomology only in
degree one, of dimension two, by
Proposition~\ref{prop:cy:specialization} at $h_0$.  Its finite torsion
and all old local data are fixed.  Choose boundary pivots there,
leaving those two free coordinates.  Their sizes and valuations may
depend on $h_0$ but not on $b$.  Apply the determinant clearing argument
of Lemma~\ref{lem:h:pivot-clearing} to the wedge of the two
bounded-integral cycles $a,2z_S$ on the remaining complex. Its allowance
for an arbitrary factor $\eta$ puts the eliminated-block factor into
the numerator alone. It yields
\begin{equation}\label{eq:cy:cleared-determinant}
 d=fu\in\mathcal O[G],\qquad f\in\mathcal O[G],\qquad
 f_0\in\Lambda,\quad f_0(0)\ne0.
\end{equation}
For each fixed $M$, the negative Laurent support of $f_x\bmod2^M$
is bounded independently of $b$ and $x$.  The clearing factor is a
polynomial of bounded degree in the remaining bounded-size matrix,
with the pivot denominators cleared by adjugates. Neither the
eliminated-block unit, the zeta class nor $d$ is required to have a uniform bound on
its negative Laurent support.

Equation~\eqref{eq:cy:rational-integrality} ensures that every $u_x$
has no negative $t$-powers.  Thus all hypotheses of
Lemma~\ref{lem:h:boolean-transfer} hold for
Equation~\eqref{eq:cy:cleared-determinant}.  In particular, for each
fixed precision and sufficiently large $b$, some nonzero $x$ satisfies
the corresponding central congruence between $u_x(0)$ and $u_0(0)$.
The dependence on $h_0$ enters the required precision and dimension,
through $v(f_0(0))$, and not the eventual valuation comparison.

In either residual case let $C$ bound the absolute error in
Equation~\eqref{eq:cy:determinant-center}.  The nonzero vertices have
$v(u_x(0))\le B+C$.  If $L(E^{(h_0)},1)=0$, the central coordinate at
zero is zero, and a congruence of precision larger than $B+C$ is
impossible.  Hence the central value is nonzero.  If
$D(h_0)>B+2C$, then
$v(u_0(0))>B+C$, giving the same contradiction.  This proves
Equation~\eqref{eq:cy:transferred-bound} with $C_E=2C$.
Only the specialization error determines this final loss; all pivot
and pole bounds determine how large a cube is needed.
\end{proof}
